\documentclass{article}
% Packages
\usepackage{amsmath}
\usepackage{amssymb}
\usepackage{enumitem}
\usepackage{hyperref}
\usepackage[margin=1in]{geometry}
\usepackage{algorithm}
\usepackage{algpseudocode} % algorithmicx is loaded by algpseudocode, but explicitly listing it as per rules
\usepackage{float} % Required for [H] float placement
\usepackage{listings} % As per rule, even if not explicitly used in this transcript
\usepackage{tikz} % As per rule, even if no network diagrams are described

% Set up document
\begin{document}

\section*{Introduction} % Use * for unnumbered sections
Welcome to this first lecture on reinforcement learning. This lecture introduces fundamental concepts that are crucial for understanding subsequent lectures. We will begin with intuitive examples, then move to a more formal presentation within the framework of Markov Decision Processes.

\section*{The Grid-World Example}
A recurring example throughout this course is the \textbf{grid-world}.
\begin{itemize}[noitemsep]
    \item It consists of a grid where a robot navigates.
    \item Grid cells can be of different types: accessible (white), forbidden (yellow), or target areas.
    \item The grid-world also has boundaries.
    \item The robot can move horizontally (left, right) and vertically (up, down) between cells, but not diagonally.
\end{itemize}
This example is chosen for its ease of understanding and illustrative utility.

\section*{The Optimal Path Problem}
The primary task in the grid-world is to find an optimal path from an initial location to a target location. This task raises several questions:
\begin{itemize}[noitemsep]
    \item How do we define a "good" or "bad" path?
    \item Intuitively, a good path avoids forbidden areas and unnecessary detours.
    \item Examples of undesirable actions include wandering aimlessly or attempting to cross grid boundaries (e.g., moving beyond the top or left edge).
\end{itemize}
Optimal policies will be discussed later in the course.

\section*{Core Concepts}

\subsection*{State}
The first core concept is \textbf{State}.
\begin{itemize}[noitemsep]
    \item A state describes the agent's status within its environment.
    \item In the grid-world, a state refers to the robot's location, e.g., $s_1, s_2, s_3, \dots, s_9$.
    \item For instance, $s_1$ corresponds to a specific $(x, y)$ coordinate on the two-dimensional grid.
    \item In more complex problems, a state might include additional information beyond location, such as speed, acceleration (for a robot), or other relevant environmental parameters.
\end{itemize}

\subsection*{State Space}
The collection of all possible states forms the \textbf{State Space}.
\begin{itemize}[noitemsep]
    \item Despite the term "space" potentially sounding complex, it is simply a set.
    \item We use the symbol $S$ to represent the state space: $S = \{s_1, s_2, \dots, s_N\}$, where $N$ is the total number of states.
    \item The range of indices (e.g., $i$ in $s_i$) can be specified or inferred from context.
\end{itemize}

\subsection*{Action}
The next fundamental concept is \textbf{Action}.
\begin{itemize}[noitemsep]
    \item In each state, an agent can take a series of actions.
    \item In our grid-world example, there are five possible actions, denoted $a_1$ to $a_5$:
        \begin{itemize}[noitemsep]
            \item $a_1$: Move Up
            \item $a_2$: Move Right
            \item $a_3$: Move Down
            \item $a_4$: Move Left
            \item $a_5$: Stay Still (represented by a circle)
        \end{itemize}
\end{itemize}

\subsection*{Action Space}
The collection of all possible actions forms the \textbf{Action Space}, represented by the set $A$.
\begin{itemize}[noitemsep]
    \item It is crucial to understand that the action space can be state-dependent.
    \item The set of available actions may vary from one state to another.
    \item Therefore, we denote it as $A(s_i)$, indicating that the action space is a function of the current state $s_i$.
\end{itemize}

\subsection*{State Transition}
The next concept is \textbf{State Transition}.
\begin{itemize}[noitemsep]
    \item This refers to the process where an agent moves from one state to another after taking an action.
    \item \textbf{Example 1 (Deterministic)}: If the agent is in state $s_1$ and takes action $a_2$ (move right), the subsequent state would be $s_2$. This transition can be represented as: $s_1 \xrightarrow{a_2} s_2$.
    \item \textbf{Example 2 (Boundary Interaction)}: If the agent is in $s_1$ and takes action $a_1$ (move up), the next state is $s_1$ again. This occurs because attempting to move up from $s_1$ encounters a boundary, causing the agent to bounce back to its original position: $s_1 \xrightarrow{a_1} s_1$.
    \item State transition defines the interaction between the agent and its environment. In a simulation or game, these interactions can be defined arbitrarily. However, in real-world scenarios, interactions are not arbitrary.
\end{itemize}

\subsubsection*{Forbidden Areas and Transitions}
Let's examine the forbidden area. Suppose an agent is in $s_5$ and takes action $a_2$ (move right). Two scenarios are possible:
\begin{enumerate}[noitemsep]
    \item \textbf{Accessible but Punished}: The forbidden area is accessible, but entering it incurs a punishment. In this case, from $s_5$, taking action $a_2$ leads to $s_6$, meaning the agent has entered the forbidden area ($s_5 \xrightarrow{a_2} s_6$).
    \item \textbf{Physically Inaccessible}: The forbidden area is physically blocked. In this scenario, from $s_5$, taking action $a_2$ would cause the agent to bounce back to $s_5$ ($s_5 \xrightarrow{a_2} s_5$).
\end{enumerate}
In this course, we adopt the first case, where forbidden areas are accessible but incur a punishment. This approach is more generalized and challenging. Allowing access to forbidden areas can lead to interesting behaviors; for example, an agent might risk entering a forbidden area to find the shortest path to the target.

\subsubsection*{State Transition Probabilities}
A tabular representation has limitations as it can only describe deterministic cases. A more general approach involves using \textbf{State Transition Probabilities}.
\begin{itemize}[noitemsep]
    \item This introduces the concept of probability into our course. Intuitively, if an agent is in $s_1$ and takes action $a_2$ (move right), the next state is $s_2$. This intuition is formalized using conditional probability, denoted $P(s' | s, a)$.
    \item \textbf{Deterministic Case}: If an agent is in $s_1$ and takes action $a_2$, the probability of moving to $s_2$ is 1.
        \begin{align*}
            P(s_2 | s_1, a_2) &= 1 \\
            P(s_i | s_1, a_2) &= 0 \quad \text{for } i \neq 2
        \end{align*}
    \item \textbf{Stochastic Case}: If wind is present, attempting to move right from $s_1$ might result in a 50\% chance of landing in $s_2$ and a 50\% chance of being pushed to $s_5$.
        \begin{align*}
            P(s_2 | s_1, a_2) &= 0.5 \\
            P(s_5 | s_1, a_2) &= 0.5
        \end{align*}
        Thus, conditional probability offers a more general framework for describing state transitions.
\end{itemize}

\subsection*{Policy}
Next, we introduce a crucial concept unique to reinforcement learning: \textbf{Policy}.
\begin{itemize}[noitemsep]
    \item A policy dictates the actions an agent should take when in a given state.
    \item Intuitively, a policy can be represented by arrows within the grid-world, where each arrow (or a circle for staying still) indicates the prescribed action for that state.
    \item A policy enables the generation of paths or trajectories. Starting from an initial location, the agent follows a specific path determined by the policy and state transitions. Upon reaching a state where the policy indicates staying still, the agent stops.
\end{itemize}

\subsubsection*{Mathematical Representation of Policy}
In reinforcement learning, the symbol $\pi$ typically denotes a policy.
\begin{itemize}[noitemsep]
    \item A policy $\pi$ is a conditional probability distribution that specifies the likelihood of taking any action in any given state. It is denoted as $P(a | s)$.
    \item The sum of probabilities for all possible actions in a given state must always equal 1: $\sum_{a \in A(s)} P(a | s) = 1$.
    \item \textbf{Deterministic Policy}: For state $s_1$, a policy might dictate that the probability of taking action $a_1$ is 0, while the probability of taking action $a_2$ is 1.
        \begin{align*}
            P(a_1 | s_1) &= 0 \\
            P(a_2 | s_1) &= 1 \\
            P(a_i | s_1) &= 0 \quad \text{for other actions } a_i
        \end{align*}
    \item \textbf{Stochastic Policy}: For instance, a policy might specify a 0.5 probability of moving right ($a_2$) and a 0.5 probability of moving down ($a_3$) from state $s_1$.
        \begin{align*}
            P(a_2 | s_1) &= 0.5 \\
            P(a_3 | s_1) &= 0.5
        \end{align*}
        (Probabilities for all other actions are 0.)
    \item This policy can also be expressed in a tabular form, where each row corresponds to a state and each column to an action. This representation is highly general, capable of describing both deterministic and stochastic policies. In programming, such policies are typically represented using arrays or matrices.
\end{itemize}

\subsubsection*{Executing Stochastic Policies}
To explain how stochastic policies are executed in practice, consider a policy with a 0.5 probability for action $a_2$ and a 0.5 probability for action $a_3$. This is implemented through random sampling.

\begin{algorithm}[H]
    \caption{Stochastic Policy Execution}
    \label{alg:stochastic_policy}
    \begin{algorithmic}[1]
        \Procedure{ChooseActionFromPolicy}{$s$, $\pi$}
            \State Let $A(s) = \{a_1, a_2, \dots, a_k\}$ be the set of available actions in state $s$.
            \State Let $P(a_i | s)$ be the probability of taking action $a_i$ given state $s$, defined by policy $\pi$.
            \State Generate a random value $x \sim U(0, 1)$.
            \State Initialize cumulative probability $C \gets 0$.
            \For{$i = 1 \dots k$}
                \State $C \gets C + P(a_i | s)$
                \If{$x < C$}
                    \State \textbf{return} $a_i$
                \EndIf
            \EndFor
        \EndProcedure
    \end{algorithmic}
\end{algorithm}
This is the fundamental principle.

\end{document}